\documentclass[12pt,a4paper]{article}
\usepackage[utf8]{inputenc}
\usepackage[T1]{fontenc}
\usepackage{geometry}
\usepackage{enumitem}
\usepackage{fancyhdr}
\usepackage{lastpage}
\usepackage{hyperref}

\geometry{margin=2.5cm}
\pagestyle{fancy}
\fancyhf{}
\rhead{Page \thepage\ of \pageref{LastPage}}
\lhead{ECO104 - Quiz (en)}
\renewcommand{\headrulewidth}{0.4pt}

\title{Quiz: ECO104}
\author{Language: en}
\date{\today}

\begin{document}

\maketitle
\thispagestyle{fancy}

\section*{Instructions}
Answer all questions by selecting the best option (A, B, C, or D). The answer key is provided at the end of this document.

\bigskip


\subsection*{Question 1}
Which of the following is NOT one of the three essential functions of money?

\begin{enumerate}[label=\Alph*., leftmargin=*]
\item Store of Value
\item Medium of Exchange
\item Unit of Account
\item Source of Energy
\end{enumerate}

\bigskip


\subsection*{Question 2}
Which of the following is necessary for a successful transaction in a barter system?

\begin{enumerate}[label=\Alph*., leftmargin=*]
\item A minimum balance requirement
\item A double coincidence of wants
\item A common unit of account
\item A central banking authority
\end{enumerate}

\bigskip


\subsection*{Question 3}
What is the historical connection between money and commodities like gold?

\begin{enumerate}[label=\Alph*., leftmargin=*]
\item Gold is unique
\item Gold is easily divisible
\item Gold makes money look more attractive
\item Gold is scarce and widely accepted as a store of value
\end{enumerate}

\bigskip


\subsection*{Question 4}
What was one advantage of metal coinage over commodity money?

\begin{enumerate}[label=\Alph*., leftmargin=*]
\item It was more easily transportable
\item It was more durable
\item It had a higher intrinsic value
\item It was not prone to counterfeiting
\end{enumerate}

\bigskip


\subsection*{Question 5}
What is one reason fiat currency leads to centralised control?

\begin{enumerate}[label=\Alph*., leftmargin=*]
\item It encourages decentralised financial systems.
\item It allows governments to control the money supply and interest rates.
\item It is backed by physical commodities like gold and silver.
\item It is resistant to inflation.
\end{enumerate}

\bigskip


\subsection*{Question 6}
How does inflation affect wage earners' purchasing power over time?

\begin{enumerate}[label=\Alph*., leftmargin=*]
\item Inflation erodes purchasing power as wages often lag behind the true rate of inflation.
\item Inflation has no effect on purchasing power.
\item Inflation only affects the purchasing power of those who invest in financial markets.
\item Inflation increases purchasing power as wages rise faster than prices.
\end{enumerate}

\bigskip


\subsection*{Question 7}
Why is it far harder to purchase a house today than in the 1980s?

\begin{enumerate}[label=\Alph*., leftmargin=*]
\item Lack of properties and suppressed interest rates make debt more enticing
\item Inflation and rental rates
\item Suppressed interest rates make debt more enticing and inflation
\item The rising cost of living and population growth
\end{enumerate}

\bigskip


\subsection*{Question 8}
What is the impact of printing more money on a currency's purchasing power?

\begin{enumerate}[label=\Alph*., leftmargin=*]
\item No impact
\item A deterioration in the purchasing power of the currency
\item More money = greater purchasing power
\item A rise in the purchasing power of the currency
\end{enumerate}

\bigskip


\subsection*{Question 9}
How are new bitcoin created and introduced into circulation?

\begin{enumerate}[label=\Alph*., leftmargin=*]
\item By receiving them as transaction fees
\item Through mining as a reward for adding new transactions to the blockchain
\item Through a centralised authority's issuance
\item By purchasing them from Bitcoin exchanges
\end{enumerate}

\bigskip


\subsection*{Question 10}
What is the significance of the quote of The Times 03/Jan/2009, Chancellor on brink of second bailout for banks, in the first Bitcoin block?

\begin{enumerate}[label=\Alph*., leftmargin=*]
\item It gives a great example of how bailouts have benefited our society
\item It signifies Satoshi’s appreciation for The Times newspaper
\item It shows us Bitcoin's ability to store messages in the blockchain
\item It highlights our financial system's ability to socialise losses
\end{enumerate}

\bigskip


\subsection*{Question 11}
What is the purpose of nodes in the Bitcoin network?

\begin{enumerate}[label=\Alph*., leftmargin=*]
\item To compete with other miners
\item To propose and implement changes to the Bitcoin codebase
\item To maintain and secure the blockchain
\item To process transactions
\end{enumerate}

\bigskip


\subsection*{Question 12}
What distinguishes Bitcoin from traditional currencies controlled by central banks or governments?

\begin{enumerate}[label=\Alph*., leftmargin=*]
\item Bitcoin requires collaboration among miners, nodes, and developers
\item All of them
\item Bitcoin allows anyone to participate in the network and perform key roles
\item Bitcoin relies on a decentralised system of consensus
\end{enumerate}

\bigskip


\subsection*{Question 13}
What advantage does Bitcoin's network offer over traditional financial systems?

\begin{enumerate}[label=\Alph*., leftmargin=*]
\item Centralised control over transactions
\item Increased inflation rates
\item Dependency on intermediaries
\item Decentralisation and security
\end{enumerate}

\bigskip


\subsection*{Question 14}
How does Bitcoin benefit individuals in war-torn or unstable countries?

\begin{enumerate}[label=\Alph*., leftmargin=*]
\item It allows access to traditional banking services
\item It provides a physical asset that can increase in value
\item It offers protection against deflation
\item It enables the storing of value in a decentralised and secure manner that has greater resistance to theft
\end{enumerate}

\bigskip


\subsection*{Question 15}
What is the role of private keys in Bitcoin transactions?

\begin{enumerate}[label=\Alph*., leftmargin=*]
\item They are shared publicly for transparency.
\item They are used to encrypt the transaction data.
\item They are used to sign and authorise transactions.
\item They are used to verify the authenticity of transactions.
\end{enumerate}

\bigskip


\subsection*{Question 16}
What is the main advantage of using non-custodial wallets for Bitcoin storage?

\begin{enumerate}[label=\Alph*., leftmargin=*]
\item Lower transaction fees
\item Easy access to Bitcoin
\item Increased security and privacy
\item No need for a password
\end{enumerate}

\bigskip


\subsection*{Question 17}
When taking self-custody, what is the first thing you should do to minimise the risk of loss?

\begin{enumerate}[label=\Alph*., leftmargin=*]
\item Disable two-factor authentication (2FA)
\item Back up your wallet recovery seed phrase
\item Take note of your recovery seed phrase and store it on your computer
\item Share your private keys with a trusted party
\end{enumerate}

\bigskip


\subsection*{Question 18}
What are the main types of stablecoins mentioned in the text?

\begin{enumerate}[label=\Alph*., leftmargin=*]
\item Stablecoins backed by banks and stablecoins backed by individuals
\item Centralised and decentralised stablecoins
\item Digital and physical stablecoins
\item Fiat-backed, commodity-backed, and algorithmic stablecoins.
\end{enumerate}

\bigskip


\subsection*{Question 19}
What was a significant limitation faced by individuals during the Greek government's 2015 bankruptcy crisis?

\begin{enumerate}[label=\Alph*., leftmargin=*]
\item The volatility of digital currencies like Bitcoin
\item Overbearing banking reform imposed by the government.
\item Limited access to digital banking services
\item A shortage of physical currency (cash)
\end{enumerate}

\bigskip


\subsection*{Question 20}
How do stablecoins, such as USDt, achieve relative stability in value?

\begin{enumerate}[label=\Alph*., leftmargin=*]
\item They rely on blockchain technology for secure transactions.
\item They are backed by a central authority or government.
\item They are pegged to a stable asset or currency.
\item They are subject to fluctuations similar to other cryptocurrencies.
\end{enumerate}

\bigskip


\subsection*{Question 21}
How do Tether’s various stablecoins contribute to financial freedom for individuals in unstable economies?

\begin{enumerate}[label=\Alph*., leftmargin=*]
\item By eliminating the need for physical cash
\item By offering higher interest rates on savings
\item By providing access to currencies with stable exchange rates
\item By enabling instant global transactions
\end{enumerate}

\bigskip


\subsection*{Question 22}
How do stablecoins help reduce the costs associated with remittances?

\begin{enumerate}[label=\Alph*., leftmargin=*]
\item By digitising fiat currencies
\item By giving individuals better exchange rates than traditional services
\item By offering lower fees than traditional services
\item By providing free money transfer services
\end{enumerate}

\bigskip


\subsection*{Question 23}
How many countries around the world restrict women from opening bank accounts in their own name?

\begin{enumerate}[label=\Alph*., leftmargin=*]
\item 72
\item 100
\item 36
\item 75
\end{enumerate}

\bigskip


\subsection*{Question 24}
What is the Halving in the context of Bitcoin?

\begin{enumerate}[label=\Alph*., leftmargin=*]
\item A mechanism that prevents Bitcoin's value from collapsing
\item A feature in Bitcoin's code that halves the mining reward every four years
\item A process of increasing the global adoption of Bitcoin
\item A term used to describe the event when Bitcoin’s price skyrockets
\end{enumerate}

\bigskip


\subsection*{Question 25}
Why is the argument that Bitcoin is worthless because it is not backed by anything flawed?

\begin{enumerate}[label=\Alph*., leftmargin=*]
\item Bitcoin is backed by gold, which gives it inherent value.
\item Bitcoin's value is determined by its scarcity and security features.
\item Value is subjective and based on people's beliefs and experiences.
\item Bitcoin's worth is based on its acceptance by governments and banks.
\end{enumerate}

\bigskip


\subsection*{Question 26}
What is the blockchain trilemma in the realm of cryptocurrencies?

\begin{enumerate}[label=\Alph*., leftmargin=*]
\item A characteristic of Bitcoin that prioritises decentralisation and security over scalability
\item The difficulty in determining the trajectory and values of cryptocurrencies
\item The limitations of Bitcoin's technology in comparison to other cryptocurrencies
\item The trade-off between scalability, security, and decentralisation in cryptocurrencies
\end{enumerate}

\bigskip


\subsection*{Question 27}
What percentage of Bitcoin's total transaction volume has been used for illicit activity?

\begin{enumerate}[label=\Alph*., leftmargin=*]
\item Less than 1\%
\item More than 5\%
\item It is unknown
\item 2-4\%
\end{enumerate}

\bigskip


\subsection*{Question 28}
Can anyone duplicate the Bitcoin code?

\begin{enumerate}[label=\Alph*., leftmargin=*]
\item No, and you can copy the value of the Bitcoin network
\item No, because the Bitcoin code is closed-source
\item Yes, but you cannot copy the value of the Bitcoin network
\item Yes, and you would copy the value of the Bitcoin network too
\end{enumerate}

\bigskip


\subsection*{Question 29}
What is one of the motivations behind Bitcoin consuming energy through mining?

\begin{enumerate}[label=\Alph*., leftmargin=*]
\item To make Bitcoin more expensive to use
\item To discourage new users from joining the network
\item To reduce its value by increasing energy costs
\item To provide a powerful defence mechanism against centralisation
\end{enumerate}

\bigskip


\subsection*{Question 30}
What is one reason why it is important to conduct thorough research before trusting a stablecoin?

\begin{enumerate}[label=\Alph*., leftmargin=*]
\item Stablecoins have no risks associated with them
\item Some stablecoins lack transparency and diligence
\item Some stablecoins are regulated by the government
\item Stablecoins are always fully backed by reserves
\end{enumerate}

\bigskip


\subsection*{Question 31}
What is the primary reason for short-term price volatility in stablecoins?

\begin{enumerate}[label=\Alph*., leftmargin=*]
\item Hacks compromising the stablecoin
\item Lack of reserves to meet withdrawal demands
\item Depegging of stablecoins
\item Low liquidity levels on exchanges
\end{enumerate}

\bigskip


\subsection*{Question 32}
In what scenario might a stablecoin trade above its pegged price?

\begin{enumerate}[label=\Alph*., leftmargin=*]
\item When the stablecoin issuer freezes funds
\item When the demand for stablecoins surpasses an exchange's capacity
\item When the stablecoin is fully backed by cash reserves
\item When there is strong selling pressure from users
\end{enumerate}

\bigskip


\subsection*{Question 33}
What is one benefit of centralised stablecoins despite their ability to freeze funds?

\begin{enumerate}[label=\Alph*., leftmargin=*]
\item They can swiftly respond to protect users in certain scenarios
\item They offer complete decentralisation and security
\item They are not subject to any regulations or interventions
\item They operate independently from law enforcement
\end{enumerate}

\bigskip


\newpage
\section*{Answer Key}

\begin{enumerate}
\item \textbf{Question 1:} D
\\ \textit{The three essential functions of money are a store of value, a medium of exchange, 
and a unit of account. A source of energy is not a function of money.}

\item \textbf{Question 2:} B
\\ \textit{In a barter system, transactions take place directly between two parties exchanging goods or 
services without the use of money. A successful transaction requires a double coincidence of 
wants, meaning that both parties must have something the other wants to trade.}

\item \textbf{Question 3:} D
\\ \textit{Given gold's globalised scarcity in nature and the significant investment required 
to mine, refine, and store it, gold has consistently traded at a premium above
other forms of money.}

\item \textbf{Question 4:} B
\\ \textit{Metallic money, particularly gold and silver, was more durable than commodity money. 
Precious metals are less prone to deterioration over time, making them ideal for
long-term storage of value. This durability also contributed to their usefulness as
a medium of exchange, a unit of account, and a store of value.}

\item \textbf{Question 5:} B
\\ \textit{Fiat currency gives governments the power to control the money supply and
interest rates, allowing them to influence economic activity and resulting
in centralised control.}

\item \textbf{Question 6:} A
\\ \textit{Inflation leads to a persistent increase in consumer prices, eroding the purchasing 
power of wage earners as their wages typically lag behind the true rate of inflation.
This makes it difficult for wage earners to achieve their long-term savings goals
and maintain their standard of living.}

\item \textbf{Question 7:} C
\\ \textit{While inflation is devaluing our currency’s purchasing power, making it 
harder to save, suppressed interest rates are making debt consumption 
more enticing. When the cost of capital is so cheap, people borrow beyond 
their means, funnelling more capital into assets, and driving up prices.}

\item \textbf{Question 8:} B
\\ \textit{When central banks decide to print money, they can’t simply create more value. 
For this newly printed money to be worth something, its value must come from the 
previous currency holders.
Think of the money supply as a pizza, and imagine it cut into four slices. 
Doubling the money supply would not be equivalent to doubling the amount of pizza. 
Instead, it would be equivalent to cutting those four slices in half to create 
eight slices. We have not gained any additional pizza. We just have more slices, 
each smaller in size. Therefore, when we print more money, we devalue the money 
that is already in existence.}

\item \textbf{Question 9:} B
\\ \textit{New bitcoin are introduced into circulation through a process called mining. 
Miners compete with one another to add new transactions to the blockchain. 
In exchange for their work, miners are rewarded with newly created bitcoin. 
This process also verifies and secures transactions on the Bitcoin network.}

\item \textbf{Question 10:} D
\\ \textit{This is a reference to a Times newspaper [article](https: //www.thetimes.co.uk/article/chancellor-alistair-darling-on-brink-of-second-bailout-for-banks-n9l382mn62h) reflecting Satoshi’s 
concerns that banks were engaging in risky behaviour, with little 
consequence to them, and that losses would be shared amongst currency holders.}

\item \textbf{Question 11:} C
\\ \textit{Nodes are essential to the integrity and security of the Bitcoin network. 
They serve as the gatekeepers of the system, ensuring that transactions 
are processed accurately and according to the rules.}

\item \textbf{Question 12:} B
\\ \textit{What distinguishes Bitcoin from traditional fiat currencies is that it's a 
collaborative effort between three essential roles and the community.
Nodes ensure the integrity and security of the network by serving as gatekeepers.
Miners are responsible for processing transactions and maintaining the blockchain.
Developers are highly skilled members of the Bitcoin community who work together 
to maintain and improve the software that powers the network.
The community, through the use of Bitcoin, is what gives it value.
Each of these roles is essential to the functionality and success of Bitcoin as a 
digital currency. By entrusting these crucial tasks to the key players, Bitcoin can 
operate as a decentralised, secure, transparent, and accountable medium of exchange for 
all users, making it significantly different from the centralised currencies we use today.}

\item \textbf{Question 13:} D
\\ \textit{Bitcoin's decentralised network provides a secure and efficient way to 
transact without the need for intermediaries or third parties. This 
reduces the risk of censorship and unauthorised changes to transactions,
as well as increases privacy.}

\item \textbf{Question 14:} D
\\ \textit{In war-torn or unstable countries, individuals may struggle to 
access traditional banking services or face the risk of theft or 
confiscation of physical assets. Bitcoin allows individuals to store value 
in a decentralised and secure manner without the need for intermediaries 
or physical assets vulnerable to risks such as inflation, theft, or confiscation.}

\item \textbf{Question 15:} C
\\ \textit{Private keys are an essential component of Bitcoin transactions.
They are used to sign and authorise transactions, providing proof 
of ownership and ensuring that only the owner can spend their bitcoin.}

\item \textbf{Question 16:} C
\\ \textit{Non-custodial wallets are where you are the sole custodian of your funds,
meaning you have complete control over your private keys. Therefore,
non-custodial wallets offer a higher level of security and privacy 
over custodial wallets since you are the only one responsible for the safety of your funds.}

\item \textbf{Question 17:} B
\\ \textit{First and foremost, back up your wallet when taking self-custody.
Hardware wallets come with a recovery seed phrase, a set of words that 
can be used to recover your private keys in case your device is lost or damaged.
Create a physical copy, such as a metal [seed plate](https: //coincodex.com/article/23147/best-metal-crypto-wallets-for-seed-phrase-storage/), of this seed phrase and 
store it in a safe place. It's important to keep this recovery seed phrase secure. 
NEVER share it with anyone.}

\item \textbf{Question 18:} D
\\ \textit{In the previous chapters, we have seen three main types of stablecoins; 
fiat-backed stablecoins, which are backed by traditional currencies 
like the US dollar; commodity-backed stablecoins, which are backed by 
commodities like gold or oil; and algorithmic stablecoins, which rely 
on a set of rules or algorithms to maintain their value. However, 
it notes that successful algorithmic stablecoins have yet to hit the market.}

\item \textbf{Question 19:} B
\\ \textit{During the Greek government's 2015 bankruptcy crisis, banks were closed, 
ATM withdrawals were limited, and a significant portion of individuals' 
bank deposits were withdrawn by the government to fund fiscal irresponsibility.
This event highlighted the risks associated with having money held within the traditional banking system.}

\item \textbf{Question 20:} C
\\ \textit{Stablecoins maintain relative stability in value by being pegged 
to a more stable asset or currency, such as the US dollar or gold.
This pegging helps mitigate the volatility often associated with other 
digital assets like Bitcoin.}

\item \textbf{Question 21:} C
\\ \textit{Tether’s various stablecoins are pegged to stable assets like the US 
dollar or gold. This stability allows individuals in unstable economies 
to move their money into a currency that is not subject to the same 
inflation rates, helping them maintain the value of their savings.}

\item \textbf{Question 22:} C
\\ \textit{Stablecoins, unlike traditional money transfer services, have significantly 
lower fees because they do not require intermediaries such as banks or money 
transfer services.}

\item \textbf{Question 23:} A
\\ \textit{In 72 countries around the world, women are not allowed to open 
bank accounts in their own names.}

\item \textbf{Question 24:} B
\\ \textit{The halving is a feature in Bitcoin's code that slashes the mining 
reward by half approximately every four years. This event impacts 
miners' revenue and exerts upward pressure on the price of bitcoin.}

\item \textbf{Question 25:} C
\\ \textit{Value is subjective and in the eye of the beholder. Bitcoin is valuable 
because people believe it to be so, just like art can be valuable despite 
being made of canvas and paint. The argument that currency must be backed by 
something of value is inaccurate, as the value of fiat currencies like the US 
dollar is based on trust and belief as well.}

\item \textbf{Question 26:} D
\\ \textit{The concept of the blockchain trilemma is the idea that cryptocurrencies 
must prioritise two out of three important characteristics: scalability, 
security, and decentralisation. Bitcoin prioritises decentralisation and security, 
which affects its scalability, but it still achieves it through alternative means.}

\item \textbf{Question 27:} A
\\ \textit{According to a report by Chainalysis in 2020, illicit activity only 
accounts for about 0.34\% of all Bitcoin transactions.}

\item \textbf{Question 28:} C
\\ \textit{While the code for Bitcoin is open source and can be copied, 
the value of Bitcoin is not derived from its code. It is derived from 
its network of users and the trust and adoption it has built up over time.}

\item \textbf{Question 29:} D
\\ \textit{One of the reasons Bitcoin consumes energy through mining is to 
provide a powerful defence mechanism against centralisation and 
undue manipulation. The distributed network of miners helps prevent 
centralised entities from co-opting the network or manipulating the rules in their favour.}

\item \textbf{Question 30:} B
\\ \textit{Not all stablecoins offer the same level of transparency. 
Therefore, it is important to conduct thorough research 
before entrusting significant wealth to any particular stablecoin.}

\item \textbf{Question 31:} D
\\ \textit{Short-term price volatility in stablecoins is often a result 
of liquidity issues, where smaller exchanges may lack the capability 
to meet customer demands, causing temporary price fluctuations.}

\item \textbf{Question 32:} B
\\ \textit{Stablecoins can trade above their pegged price when the 
demand for stablecoins exceeds an exchange's capacity to 
meet customer demands, typically on smaller exchanges.}

\item \textbf{Question 33:} A
\\ \textit{Centralised stablecoins have the ability to intervene 
when it is in the best interest of users and the overall stability 
of the system, as demonstrated by examples where frozen funds were 
used to protect consumers and prevent hackers from profiting.}


\end{enumerate}

\end{document}
